\hypertarget{appendix-3cb49709605a}{}
\section*{Member order and delegating constructors}
\section{Member order and delegating constructors}\label{member-order-and-delegating-constructors}

\textbf{Curriculum depth:} Introduced · \textbf{First seen:} C009 · \textbf{Latest development:} C009

\subsection{Definition and mechanism}\label{definition-and-mechanism}

Members initialize in declaration order, not initializer-list spelling order. A provider must precede a dependent member so reverse destruction cleans the dependent first. Delegation uses one target constructor as the sole initializer; the target completes before the delegating body.

\subsection{Prerequisites}\label{prerequisites}

\begin{itemize}
\tightlist
\item
  \hyperref[appendix-7f8ec83d0886]{Constructors and invariant establishment} --- Initialization form and structural order determine how the constructor establishes members.
\end{itemize}

\subsection{Used by}\label{used-by}

\begin{itemize}
\tightlist
\item
  \hyperref[appendix-9bfe377cef7a]{Destruction and partial construction}
\end{itemize}

\subsection{Recall and next depth}\label{recall-and-next-depth}

Put a device and its dependent buffer in lifetime order; explain a delegating validation path. Pending: virtual bases and inheritance construction. Written coverage does not demonstrate personal mastery.

\subsection{Source and navigation}\label{source-and-navigation}

\href{Continuous\%20Notes/Computing\%20Foundations\%20-\%20From\%20Bits\%20to\%20C++\%20and\%20Systems\%20APIs\#Construction\%20order\%20records\%20lifetime\%20dependencies}{Teaching mechanism and examples} · \hyperref[appendix-fd090d982a15]{Exact coverage} · \href{}{Dependency Map (02-Foundations-II, appendix-4ad47bdd269a)}
